\documentclass[11pt]{article}
\usepackage[a4paper,margin=27mm]{geometry}
\usepackage{amsmath,amssymb,mathtools}
\usepackage{fontspec,unicode-math}
\setmainfont{Latin Modern Roman}
\setmathfont{Latin Modern Math}
\usepackage{xurl}
\usepackage[hidelinks]{hyperref}
\newcommand{\sref}[2]{\hyperlink{source:#1:#2}{[S#2]}}
\newcommand{\cataloguescope}[1]{\paragraph{Mathematical question.}#1}
\hypersetup{pdftitle={Problem statement}}
\begin{document}
\section{Power-of-two periodicity conjecture}\label{statement}

\subsection{Definitions and mathematical statement}

\paragraph{Shared notation.}$\mathbb N=\{1,2,\ldots\}$, and $\mathbb Z,\mathbb Q,\mathbb R,\mathbb C$ denote the integers, rationals, reals and complex numbers. Any other conventions are specified locally.


Define the Euler up/down numbers $E_n$ for integers $n\ge0$ by
\[
\sum_{n\ge0} E_n\frac{z^n}{n!}=\sec z+\tan z.
\]
For $q\in\mathbb N$ this sequence modulo $q$ is eventually periodic. Its minimal preperiod is
\[
s(q)=\min\{s\in\mathbb Z_{\ge0}:\exists d\in\mathbb N\ \forall n\ge s,\ E_{n+d}\equiv E_n\pmod q\},
\]
and its minimal eventual period is
\[
d(q)=\min\{d\in\mathbb N:\forall n\ge s(q),\ E_{n+d}\equiv E_n\pmod q\}.
\]
The prime-level periods are $d(p)=p-1$ for $p\equiv1\pmod4$ and $d(p)=2(p-1)$ for $p\equiv3\pmod4$.

For $n\ge1$ and $1\le i\le n$, let $e_{n,i}$ count permutations $\sigma$ of $\{1,\ldots,n\}$ such that $\sigma(n)=i$ and $\sigma(j-1)<\sigma(j)$ for even $j$, while $\sigma(j-1)>\sigma(j)$ for odd $j$. In particular $e_{1,1}=1$; equivalently, for $n\ge2$,
\[
e_{n,i}=\begin{cases}\sum_{j<i}e_{n-1,j},&n\text{ even},\\\sum_{j\ge i}e_{n-1,j},&n\text{ odd}.\end{cases}
\]
In these sums $1\le j\le n-1$. Write $v_2(m)$ for the exponent of $2$ in a positive integer $m$, with $v_2(0)=+\infty$. The $i$th diagonal of the valuation triangle consists of $v_2(e_{n,i})$ for $n\ge i$. Set
\[
m_i=\min_{n\ge i}v_2(e_{n,i}),\qquad
u_k=\sup\bigl(\{i\in\mathbb N:m_i<k\}\cup\{0\}\bigr)\quad(k\ge1),
\]
where the supremum is allowed to be $+\infty$. When finite this is the maximum used in Arnold's sequence. Thus a proposed equality with a finite integer also asserts the required finiteness. The recorded initial values of $u_k$ are $2,4,4,4,8,8,8,8,10,12,12,16,\ldots$.
\paragraph{Statement recorded in the supplied source.}
For every integer $k\ge1$, is $s(2^k)=u_k$? Also, is $d(2^k)=2^k$ for every $k\ge1$ with $k\ne2$, with the exceptional value $d(4)=2$? \sref{statement}{2}
\subsection{Short English statement}
Are the preperiods modulo powers of two given by Arnold's valuation sequence, and are the periods the corresponding powers of two except for modulus four?
\subsection{Sources}
\begin{description}
\item[\hypertarget{source:statement:1}{S1}] ULAM.AI and the UnsolvedMath Contributors, UnsolvedMath: A Curated Collection of Open Mathematics Problems, catalogue entry. CC BY 4.0. \url{https://huggingface.co/datasets/ulamai/UnsolvedMath}.
\item[\hypertarget{source:statement:2}{S2}] Sanjay Ramassamy, \emph{Modular periodicity of the Euler numbers and a sequence by Arnold}, Arnold Mathematical Journal (2018) \S3, equations (5)--(7), Conjecture 2, pp. 3--5. \url{https://amj.math.stonybrook.edu/html-articles/Files-2015-2024/18-79/index.html}.
\end{description}
\label{end:statement}
\end{document}
